\subsection{CTR: comprobación de la tupla, exposición y fuerza concreta}
\label{sec:ctr-verified}
Las declaraciones indicadas aquí pertenecen al módulo
\nolinkurl{CTRProfile.lean}, espacio de nombres \nolinkurl{CTRProfile}. El modelo
usa palabras de 128 bits, claves formadas por dos palabras y el contador
completo módulo \(2^{128}\). El argumento exacto de las propiedades
estructurales usa las dos leyes de inversión del cifrado de bloque para cada
clave. La hipótesis de dureza de AES interviene en los pasos computacionales.

\paragraph{Paso 1. Del procedimiento a las tuplas.}
Una transcripción itera el incremento y el cifrado del contador, acumula
tres bloques para cada \emph{Update}, conserva sus 384 bits, aplica el XOR
y divide el resultado en clave y contador. Una segunda transcripción expresa
la tupla de la sección~\ref{sec:ctr} por sus tres cifrados explícitos.
\proofat{CTRProfile.update_refines}{update\_\allowbreak{}refines}, \proofat{CTRProfile.generate_refines}{generate\_\allowbreak{}refines} y
\proofat{CTRProfile.checked_refines}{checked\_\allowbreak{}refines} establecen igualdad de la salida, la clave,
el contador criptográfico y el contador administrativo en todas las entradas
admitidas, con la misma respuesta si hace falta resembrar. Incluyen entrada
ausente, entrada de 384 bits presente y, como caso particular, una entrada
presente de ceros. El primer caso omite la actualización previa y el último
la ejecuta. La petición de un bloque satisface DRG.3.2. En el modelo,
\(\mathrm{inc}(2^{128}-1)=0\) y los contadores consecutivos que se
emplean en la prueba siguen siendo distintos
\cite[§§10.2.1.2 y 10.2.1.5.1]{nist90a}.

Las dos 4-tuplas completas son
\[
 (SM,PS,\mathcal S,\phi_{\rm seed}),\qquad
 (SM,PS,\mathcal S,\phi_{\rm reseed}),
\]
con
\[
 SM=\bits^{384},\quad
 PS=\bigcup_{j=0}^{384}\bits^j,\quad
 \mathcal S=\bits^{256}\times\bits^{128}.
\]
Para \(p\in PS\), \(\bar p=p\|0^{384-|p|}\). Tomando cero como el par
de estado nulo, la primera función es
\(\phi_{\rm seed}(x,p)=\mathrm{Upd}(0,x\oplus\bar p)\).
La segunda es
\(\phi_{\rm reseed}(s,x,p)=\mathrm{Upd}(s,x\oplus\bar p)\).
Ambas reinician el contador administrativo en 1. Lean recibe la cadena ya
rellenada, en las tres palabras de 128 bits en su orden de transmisión.
\proofat{CTRProfile.instantiate_bijection}{instantiate\_\allowbreak{}bijection} y \proofat{CTRProfile.reseed_bijection}{reseed\_\allowbreak{}bijection} demuestran
que cada estado tiene una única semilla preimagen para el contexto fijado.
En la siembra noDF de 90A no se usa nonce. El contador público y la
personalización tienen crédito de entropía cero. La traducción preserva
exactamente la distribución conjunta de la semilla salvo reetiquetado.
Para .1 y .4, la evidencia de fuente debe respaldar el origen y la entropía
de la parte efectiva tras fijar el contexto. La exigencia NIST del mecanismo
es una entrada de 384 bits de \emph{full entropy} en siembra y resiembra
\cite[§§143, 147 y 205--207]{ais}\cite[§§10.2.1.3.1 y 10.2.1.4.1]{nist90a}.

Puede precisarse la obligación de minentropía para la opción TRNG de
DRG.3.4. Sea \(E\) el contexto clásico que comprende el estado antiguo y
la entrada de personalización o resiembra. Por la biyección anterior,
\(\hm(K^+,V^+\mid E)=\hm(SM\mid E)\). Incluso concediendo \(V^+\)
al observador, la definición de adivinación sobre masas conjuntas da
\[
\begin{aligned}
  \pg(K^+\mid V^+,E)
    &=\sum_{e,v}\max_k P_{E,V^+,K^+}(e,v,k)\\
    &\le 2^{128}\sum_e\max_{v,k}P_{E,V^+,K^+}(e,v,k).
\end{aligned}
\]
\proofat{CTRProfile.projectedGuessBound}{projectedGuessBound} verifica la
desigualdad de masas finitas. Una cota conjunta de 368 bits en el material basta así para acreditar 240
en la coordenada de clave bajo la interpretación del estado efectivo de §102.
A modo de cálculo, tres bloques físicos dedicados de 128 bits que acrediten
\(0.98\) de minentropía por bit \emph{puntualmente} dado el historial de cada
bloque y la vista requerida aportan al menos \(376.32\) bits conjuntos;
la clave conserva al menos \(248.32\) bits de minentropía condicional
promediada. Este cálculo da una condición suficiente para la propiedad
de semilla AIS. El origen PTG.2, el soporte de estados efectivos y el
requisito de 384 bits de \emph{full entropy} de 90A sin DF tienen
justificaciones propias \cite[§§102,207,1074--1075]{ais}.

\paragraph{Paso 2. ITS y las tres vistas del estado.}
Para el perfil físico \(D=0^{256}\|A\), sean \(k'=G_K(s)\),
\(v'=G_V(s)\oplus A\) y \(E_{k'}\) el cifrado AES-256. Se verifica
\(Y=E_{k'}(\mathrm{inc}(v'))\), y
\(A=G_V(s)\oplus \mathrm{dec}(E_{k'}^{-1}(Y))\), donde
\(\mathrm{dec}=\mathrm{inc}^{-1}\). Las declaraciones
\proofat{CTRProfile.physical_output}{physical\_\allowbreak{}output}, \proofat{CTRProfile.recover_encode}{recover\_\allowbreak{}encode} y
\proofat{CTRProfile.encode_recover}{encode\_\allowbreak{}recover} dan la permutación para todo \(s\).
Para estado previo fijo y contexto lateral común se aplica la conservación
de la probabilidad óptima de adivinar y de la distancia a uniforme por una
permutación. Esto también sustenta la composición temporal del teorema de ITS.

Al leer \emph{sólo} el estado posterior, las mitades de la nueva clave
permiten calcular el indicador del suceso \(B=\{Y=C_1\lor Y=C_2\}\), con
\[
       C_1=E_{k'}(\mathrm{inc}^2(v')),\qquad
       C_2=E_{k'}(\mathrm{inc}^3(v')) .
\]
Las tres entradas de AES son distintas. Luego \(C_1\ne C_2\) y
\(\Pr_{\rm real}(B)=0\). Frente a una salida ideal uniforme e independiente
del estado posterior, \(\Pr_{\rm ideal}(B)=2/2^{128}\). Por la
caracterización de la distancia total en variación por sucesos,
\[
   \dist(P_{Y,S^+},U_{128}\otimes P_{S^+})\ \ge\ 2^{-127}.
\]
Es una cota inferior para el cifrado real incluso con entrada física
perfectamente uniforme y secreta. \proofat{CTRProfile.one_snapshot_excludes_two}{one\_\allowbreak{}snapshot\_\allowbreak{}excludes\_\allowbreak{}two}
comprueba la exclusión y la distinción, y
\proofat{CTRProfile.uniform_event_has_two}{uniform\_\allowbreak{}event\_\allowbreak{}has\_\allowbreak{}two} cuenta los casos favorables ideales.

Con \emph{ambos} estados adyacentes conocidos, bastan \(s\) y \(C_1\).
Se calcula \(k'=G_K(s)\) y
\[
              Y=E_{k'}\bigl(\mathrm{dec}(E_{k'}^{-1}(C_1))\bigr).
\]
Así, \(p_{\rm guess}(Y\mid s,C_1)=1\) y la distancia conjunta a una salida
uniforme independiente es exactamente \(1-2^{-128}\). La fórmula inversa
anterior recupera también \(A\). Los lemas son
\proofat{CTRProfile.two_snapshots_recover_output}{two\_\allowbreak{}snapshots\_\allowbreak{}recover\_\allowbreak{}output} y
\proofat{CTRProfile.two_snapshots_recover_input}{two\_\allowbreak{}snapshots\_\allowbreak{}recover\_\allowbreak{}input}.
El mismo cálculo prueba una obligación positiva de PTG.3.2: para cada
estado antiguo, dos bloques físicos distintos producen primeras mitades de
clave posterior distintas (\proofat{CTRProfile.state_influenced}{state\_\allowbreak{}influenced}, AIS §381).
El §213 de AIS analiza la divulgación del estado posterior y, cuando
corresponde, de la última entrada adicional. Su objetivo para .7 es
computacional y la vista de dos capturas añade el estado anterior.

\paragraph{Paso 3. Estado efectivo y predicción.}
Fijado un bloque autónomo conocido \(y\), a cada clave candidata \(k\)
corresponde exactamente un estado anterior
\[
            s_k=(k,\mathrm{dec}(E_k^{-1}(y))).
\]
Todos los \(s_k\) son alcanzables mediante siembra. \proofat{CTRProfile.single_output_fiber}{single\_\allowbreak{}output\_\allowbreak{}fiber}
y \proofat{CTRProfile.candidate_injective}{candidate\_\allowbreak{}injective} verifican que la fibra de un bloque tiene
exactamente un representante por clave, incluso si los siguientes bloques
pertenecen a peticiones distintas. Constituye la parte algebraica del
razonamiento del §102. Para ver por qué hace falta una premisa computacional,
\proofat{CTRProfile.fiber_size_is_insufficient}{fiber\_\allowbreak{}size\_\allowbreak{}is\_\allowbreak{}insufficient} comprueba un control negativo:
con la familia de permutaciones identidad, la fibra conserva sus
\(2^{256}\) elementos, pero el siguiente bloque autónomo vale
\(\mathrm{inc}^4(y)\) para toda clave. AIS §§91--95 y 115 requieren
resistencia de los estados efectivos también tras conocer secuencias.

Una forma de ligar el tamaño efectivo a un predictor equivalente consiste
en fijar dos bloques autónomos consecutivos futuros de 128 bits cada uno.
Para cada historial permitido, supóngase que un código de estado eficazmente
utilizable tiene a lo sumo \(L\) valores y basta para calcular esos bloques.
Si se sortea un código equiprobable con su decodificador, al menos uno acierta
para cada estado cubierto. Su probabilidad de éxito es por tanto
\(\ge 1/L\), también al promediar el historial. Las declaraciones
\proofat{CTRProfile.codebook_success}{codebook\_\allowbreak{}success} y \proofat{CTRProfile.predictor_space_lower_bound}{predictor\_\allowbreak{}space\_\allowbreak{}lower\_\allowbreak{}bound}
verifican el conteo y la implicación
\[
         p_{256}(T)\le2^{-246}\quad\Longrightarrow\quad L\ge2^{246}.
\]
Aquí \(T\) debe incluir el coste de muestrear y decodificar el código. Es
un criterio suficiente para excluir representaciones predictivas más pequeñas
en ese presupuesto. La cota \emph{aditiva} de la
sección~\ref{sec:ctr-prp} para \(N=1024\) incluye ya
\(12288\,2^{-128}\). Por sí sola no alcanza ese umbral. Para emplear la
cota unilateral de predicción hay que instanciar la calidad de semilla,
el número de tramos y la ventaja PRP admisible para el tiempo de ataque.
BSI asigna 256 bits en la discusión del §102 y exige al solicitante la
prueba del perfil en §1007. Ninguna enumeración de los 384 bits almacenados
satisface por sí misma este requisito \cite[§§91--103, 115 y 1007]{ais}.

\paragraph{Paso 4. Alcance del argumento computacional.}
Para .5 y .6 se especifican las posiciones protegidas, el prefijo o sufijo
divulgado y la longitud \(n\) del objetivo. Para .7 se añade el estado
terminal y, cuando procede, el último dato adicional. Para .8, los datos
adicionales pueden quedar fijados por el historial público ya observado,
incluida la elección de entrada ausente. Si en una simulación una petición
futura copiase al historial un bloque objetivo que ya hubiese sido mostrado
al adversario, esa elección describiría un objetivo previamente divulgado.
Los desafíos de predicción emplean posiciones ocultas y una política de
entradas que sólo recibe la vista admitida. Cada tramo usa puntos distintos
bajo una única clave y un XOR fijado antes de procesar sus respuestas.
\proofat{CTRProfile.mask_cancels}{mask\_\allowbreak{}cancels} y \proofat{CTRProfile.exposed_triple}{exposed\_\allowbreak{}triple} verifican el tratamiento
de la máscara y del último estado, para los 384 bits completos del dato y
también para la rama autónoma. A partir de ahí intervienen la hipótesis
PRP de AES, el límite de época y la justificación de la semilla de la
sección~\ref{sec:ctr-prp}. La predicción de un bloque y la predicción de
\(n>128\) bits tienen objetivos de trabajo distintos bajo §214.

Para .9, una preimagen válida de la transición, con máscara pública fijada,
debe reproducir los tres cifrados de \proofat{CTRProfile.exposed_triple}{exposed\_\allowbreak{}triple}.
El argumento de inversión debe aplicarse con la distribución concreta de
esa terna y la vista autorizada. Para .10 se adjunta evidencia de la salida
y su justificación conforme a §§157--162 y 200. Un vector funcional confirma
la transcripción para ese ejemplo y se registra aparte de la evidencia
estadística del producto. El informe sobre navegadores atribuye plausibilidad
a los mecanismos evaluados y no evalúa nuestra fuente física ni el perfil
de entrada en \(V\) \cite[§§3.2.4 y 4.2.4]{bsi-browsers}.
